\documentclass{article}
\usepackage[left=1in, right=1in, bottom = 1in, top = 1in]{geometry}
\usepackage{graphicx} % Required for inserting images
\usepackage{bbm}
\usepackage{amsmath}
\usepackage[english]{babel}
\usepackage[utf8]{inputenc}
\usepackage[authoryear]{natbib}
%\usepackage[margin=1in]{geometry}
\usepackage{setspace}

%\setlength{\parskip}{1em}
%\onehalfspacing
%\linespread{1.45}

% Math Packages
\usepackage{amsmath}
%\numberwithin{equation}{section}
\usepackage{amsfonts}
\usepackage{amsthm}
\usepackage{amssymb}
\usepackage{mathrsfs}
\usepackage{mathtools}
%\usepackage{bbm}
\usepackage{amstext}
\usepackage{comment}
\usepackage{epsfig}
\usepackage{latexsym}
\usepackage{epstopdf}
\usepackage{bm}
\usepackage{algorithm}
\usepackage{algpseudocode}
%\usepackage[capposition=top]{floatrow}


% For drawing graph
\usepackage{tikz}
\usetikzlibrary{positioning,chains,fit,shapes,calc}
\definecolor{myblue}{RGB}{80,80,160}
\definecolor{mygreen}{RGB}{80,160,80}

\usepackage{url}
\usepackage[colorlinks = true, linkcolor = blue]{hyperref}
\hypersetup{
        urlcolor = blue,
	colorlinks = true,
	linkcolor = blue,
	citecolor = blue,
}
\usepackage{paralist}
\usepackage{enumitem}
\setlist[enumerate]{topsep=0pt,noitemsep}

% For Tables
\usepackage{multirow}
\usepackage{multicol}
\usepackage{babel}
\usepackage{array}
\usepackage{rotating}
\usepackage{longtable}
\usepackage{float}
\usepackage{booktabs}
\usepackage{lscape}
\usepackage{caption}
\usepackage{subcaption}
\usepackage{afterpage}
% \usepackage{showkeys}
\usepackage{authblk}

% Theorem Environment
\newtheorem{theorem}{Theorem}[section]
\newtheorem{remark}{Remark}[section]
\newtheorem{assumption}{Assumption}
\newtheorem{lemma}{Lemma}[section]
\newtheorem{notation}{Notation}
\newtheorem{example}{Example}
\newtheorem{definition}{Definition}[section]
\newtheorem{proposition}{Proposition}[section]
\newtheorem{corollary}{Corollary}[section]
\newtheorem{prediction}{Prediction}
\newtheorem{regression}{Regression}

% Custom environment for continuance
\newenvironment{continuance}[1]
{%
  \newcommand{\continuanceref}{\ref{#1}}% Define reference to original example
  \continuancex % Begin the non-numbered theorem-like environment
}
{\endcontinuancex}

% Define a non-numbered continuation environment
\newtheorem*{continuancex}{Example \continuanceref \ (continued)}

% \makeatletter
% \renewenvironment{proof}[1][\proofname]{%
%   \par\pushQED{\qed}\normalfont%
%   \topsep6\p@\@plus6\p@\relax
%   \trivlist\item[\hskip\labelsep\bfseries#1\@addpunct{.}]%
%   \ignorespaces
% }{%
%   \popQED\endtrivlist\@endpefalse
% }
% \makeatother
%\numberwithin{theorem}{section}
%\numberwithin{corollary}{section}
%\numberwithin{proposition}{section}
%\numberwithin{lemma}{section}
%\numberwithin{assumption}{section}

% Formatting 
\linespread{1.4}
% \mathtoolsset{showonlyrefs=true}

\allowdisplaybreaks


% Define math operators
\DeclareMathOperator*{\argmax}{argmax}
\DeclareMathOperator*{\argmin}{argmin}
\DeclareMathOperator{\tr}{tr}
\DeclareMathOperator{\vect}{vec}
\DeclareMathOperator{\Rank}{Rank}

% Define shortcut
\newcommand{\E}{\mathbb{E}}
\newcommand{\B}{B}
\newcommand{\Var}{\mathrm{Var}}
\newcommand{\Cov}{\mathrm{Cov}}
\newcommand{\bbm}[1]{\mathbbm{#1}}
\newcommand{\bb}[1]{\mathbb{#1}}
\newcommand{\floor}[1]{\lfloor #1 \rfloor}
\newcommand{\sq}[1]{\textquoteleft #1 \textquoteright}
\newcommand{\dq}[1]{\textquotedblleft #1 \textquotedblright}
\newcommand{\red}{\color{red}}
\newcommand{\norm}[1]{\left\lVert#1\right\rVert}
\def\T{{\textsc{t}}}

\newcommand{\pr}{\textup{pr}}
\newcommand{\cov}{\textup{cov}}
\newcommand{\var}{\textup{var}}
\newcommand{\avar}{\textup{avar}}
\newcommand{\acov}{\textup{acov}}
\newcommand{\kprime}{k^{\prime}}
\newcommand{\adj}{\textup{adj}}

\newcommand{\bOneZ}{\boldsymbol{1}_{m}}
\newcommand{\bZeroZ}{\boldsymbol{0}_{m}}
\newcommand{\bOneA}{\boldsymbol{1}_{n}}
\newcommand{\bZeroA}{\boldsymbol{0}_{n}}
\newcommand{\Yt}[1]{Y_{#1}(\boldsymbol{1}_{m})}
\newcommand{\Yc}[1]{Y_{#1}(\boldsymbol{0}_{m})}

\newcommand{\bY}{\boldsymbol{Y}}
\newcommand{\bz}{\boldsymbol{z}}
\newcommand{\bZ}{\boldsymbol{Z}}
\newcommand{\ba}{\boldsymbol{a}}
\newcommand{\bA}{\boldsymbol{A}}
\newcommand{\bW}{\boldsymbol{W}}
\newcommand{\bV}{\boldsymbol{V}}
\newcommand{\bX}{\boldsymbol{X}}
\newcommand{\Davg}{D_{\textup{avg}}}
\newcommand{\bOmega}{\boldsymbol{\Omega}}

\newcommand{\cN}{\mathcal{N}}
\newcommand{\cF}{\mathscr{F}}
\newcommand{\cT}{\mathcal{T}}
\newcommand{\cS}{\mathcal{S}}

\newcommand{\allt}{\left(\prod_{k\in \mathcal{G}_{+i}} Z_k\right)}
\newcommand{\alltj}{\left(\prod_{k\in \mathcal{G}_{+j}} Z_k\right)}
\newcommand{\allc}{\left\{\prod_{k\in \mathcal{G}_{+i}} (1-Z_k)\right\}}
\newcommand{\allcj}{\left\{\prod_{k\in \mathcal{G}_{+j}} (1-Z_k)\right\}}

\newcommand{\indicator}{\mathbbm{1}}
\newcommand{\plim}{\textup{plim}}

\newcommand{\bG}{\boldsymbol{G}}
\newcommand{\G}{G}
\newcommand{\cG}{\mathcal{G}}
\newcommand{\cGijcap}{\cG_{+i} \cap \cG_{+j}}
\newcommand{\cGijcup}{\cG_{+i} \cup \cG_{+j}}

\newcommand{\Smax}{{\max_{1\leq i\leq n}\G_{+i}}}
\newcommand{\Dmax}{{\max_{1\leq k\leq m}\G_{k+}}}
\newcommand{\Smaxsupp}{{\bar{D}_\textup{O}}}
\newcommand{\Dmaxsupp}{{\bar{D}_\textup{I}}}

\newcommand{\sumi}{\sum_{i=1}^n}
\newcommand{\sumj}{\sum_{j=1}^n}
\newcommand{\sumk}{\sum_{k=1}^m}
\newcommand{\sumij}{\sum_{i=1}^{n} \sum_{j=1}^{n}}
\newcommand{\sumu}{\sum_{u=1}^n}
\newcommand{\sumv}{\sum_{v=1}^n}
\newcommand{\maxi}{\max_{1\leq i\leq n}}

\newcommand{\de}{\textsc{de}}
\newcommand{\ie}{\textsc{ie}}
\newcommand{\horthom}{\textup{ht}}
\newcommand{\haj}{\textup{haj}}

\newcommand{\sz}[1]{{\textcolor{olive}{(Sizhu: #1)}}}

\makeatletter
\newcommand{\ostar}{\mathbin{\mathpalette\make@circled\star}}
\newcommand{\make@circled}[2]{%
  \ooalign{$\m@th#1\smallbigcirc{#1}$\cr\hidewidth$\m@th#1#2$\hidewidth\cr}%
}
\newcommand{\smallbigcirc}[1]{%
  \vcenter{\hbox{\scalebox{0.77778}{$\m@th#1\bigcirc$}}}%
}
\makeatother


%%%%%%%%%%%%%%%%%%%%%%%%%%%%%%%%%%%%%%%%%

\title{Estimating within-cluster and between-cluster spillover effects in randomized saturation designs}

\author{ 
Sizhu Lu\thanks{Department of Statistics, University of California, Berkeley, CA 94720, USA. Email: \href{mailto:sizhu_lu@berkeley.edu}{sizhu\_lu@berkeley.edu}}, 
Lei Shi\thanks{Division of Biostatistics, University of California, Berkeley, CA 94720, USA. Email: \href{mailto:leishi@berkeley.edu}{leishi@berkeley.edu}}, 
Peng Ding\thanks{Department of Statistics, University of California, Berkeley, CA 94720, USA. Email: \href{mailto:pengdingpku@berkeley.edu}{pengdingpku@berkeley.edu}}
}

% \author[1]{Sizhu Lu\thanks{The first two authors contributed equally to this work.}}
% \author[2]{Lei Shi$^{*}$}
% \author[1]{Peng Ding}

% \affil[1]{Department of Statistics, University of California, Berkeley}
% \affil[2]{Division of Biostatistics, University of California, Berkeley}

\date{\today}
\begin{document}

\maketitle
\begin{abstract}
    Randomized saturation designs are two-stage experiments: they first randomly assign treatment probabilities over the clusters and then randomly assign the treatment to the units within the clusters. The existing literature on randomized saturation designs focuses on estimating within-cluster spillover effects by assuming away between-cluster spillover effects. However, the units may interact across clusters in practice. A leading example is that some units are geographically close to each other, so spillover effects arise across clusters. Based on the potential outcomes framework, we formulate the causal inference problem of estimating within-cluster and between-cluster spillover effects in randomized saturation designs. We clarify the causal estimands and establish rigorous statistics theory for inference. We also apply our theory to analyze a recent randomized saturation design of cash transfer on household expenditure in Kenya.
\end{abstract}
% {\small\hspace{2em} \textbf{Key Words: } }


\section{Research question}
Randomized saturation designs are widely used to study causal effects in the presence of network interference \citep{hudgens2008toward}. Most existing studies focus on within-cluster interference and assume no between-cluster spillover effects, which is often an unrealistic assumption in applied settings. A leading example is the study by \citet{egger2022general}, which implemented a randomized experiment to evaluate the impact of a cash transfer program in rural Kenya. The study used a randomized saturation design: in the first stage, sublocations were randomly assigned to either high- or low-saturation groups. In the second stage, within each high-saturation sublocation, two-thirds of villages were randomly assigned to treatment, while in low-saturation sublocations, one-third of villages were treated. All eligible households in treated villages received transfers. In this setting, there are two plausible sources of interference: (i) within-sublocation interference, where villages in the same sublocation affect each other’s outcomes, and (ii) between-sublocation interference, where villages in different sublocations, but located near each other, may also exert influence.

To formalize these ideas, let $k(i)$ denote the sublocation to which village $i$ belongs, and let $\cG_{+i}$ represent the set of villages that are geographically close to village $i$. We define the exposure mapping as $d_i(\bA) = (A_{i}, S_{i}, H_{i})$, where $S_i$ and $H_i$ are summaries of within-sublocation exposure and between-sublocation exposure \citep{aronow2017estimating}:
\begin{eqnarray*}
    S_{i} = f_i\{(A_{j}: k(j)=k(i)\}, \quad 
    H_{i} = g_i\{A_j: k(j)\neq k(i), j\in\cG_{+i}\}. 
\end{eqnarray*}
For example, we can use the following exposure mapping:  
\begin{eqnarray*}
    S_{i} &=& \indicator\left\{\frac{\sum_{j\neq i}\indicator\{k(j)=k(i)\} A_{j}}{\sum_{j\neq i}\indicator\{k(j)=k(i)\}} > \frac{1}{2}\right\},  \\
    H_{i} &=& \indicator \left\{\frac{\sumj \indicator\{k(j)\neq k(i), j\in\cG_{+i}\}A_{j}}{\sumj \indicator\{k(j)\neq k(i), j\in\cG_{+i}\}} > \frac{1}{2}\right\}.
\end{eqnarray*}


\section{Causal parameters and estimation}
\label{sec::parameters}
\subsection{Direct effect}
We first consider the direct effect of treatment. Define
\begin{eqnarray*}
    \de(s,h) &=& n^{-1}\sumi Y_{i}(1,s,h) - n^{-1}\sumi Y_{i}(0,s,h), \\
    \de &=& n^{-1}\sumi E_{\bA_{(-i)}}\left\{Y_{i}(1,S_{i},H_{i})\right\} - n^{-1}\sumi E_{\bA_{(-i)}}\left\{Y_{i}(0,S_{i},H_{i})\right\}.
\end{eqnarray*}
where $E_{\bA_{(-i)}}(\cdot)$ denotes the expectation with respect to the distribution of $\bA_{(-i)}$, the vector of treatment assignments for all villages other than $i$. 

The quantity $\de(s,h)$ represents the \textit{controlled direct effect} of treatment $A_{i}$ on the outcome, while holding the exposure variables $(S_{i},H_{i})$ fixed at values $(s,h)$. In contrast, $\de$ represents the \textit{marginal direct effect} of treatment $A_{i}$, marginalizing over the distribution of treatments for all other villages. Therefore, by definition, $\de$ depends on the treatment assignment mechanism.

\subsection{Indirect effect}
We consider two types of indirect effects corresponding to the two sources of interference: within-cluster and between-cluster spillovers.

For within-cluster indirect effects, we define the \textit{within-cluster controlled indirect effect} of $S_{i}$ on the outcome as
\begin{eqnarray*}
    \ie^{\textup{w}}(s,s^{\prime},h) &=& n^{-1}\sumi Y_{i}(0,s,h) - n^{-1}\sumi Y_{i}(0,s^{\prime},h),
\end{eqnarray*}
and the \textit{marginal within-cluster indirect effect} of $S_{i}$ on the outcome as
\begin{eqnarray*}
    \ie^{\textup{w}}(s,s^{\prime}) &=& n^{-1}\sumi E_{\bA_{(-i)}\mid S_{i} = s}\left\{Y_{i}(0,s,H_{i})\right\} - n^{-1}\sumi E_{\bA_{(-i)}\mid S_{i} = s^{\prime}}\left\{Y_{i}(0,s^{\prime},H_{i})\right\},
\end{eqnarray*}
where $E_{\bA_{(-i)}\mid S_{i} = s}(\cdot)$ denotes expectation with respect to the distribution of $\bA_{(-i)}$ conditional on $S_{i}=s$. The quantity $\ie^{\textup{w}}(s,s^{\prime},h)$ captures the effect of changing the proportion of treated villages in the same sublocation (from $s$ to $s^{\prime}$) on the outcome of village $i$, while holding its own treatment fixed at zero and the between-cluster exposure $H_{i}$ at level $h$. The $\ie^{\textup{w}}(s,s^{\prime})$ is a marginal version of this effect, marginalized over the conditional distribution of other villages' treatments given $S_{i}=s$ and $S_{i}=s^{\prime}$, respectively. As with the marginal direct effect, this definition depends on the treatment assignment mechanism.

We similarly define between-cluster indirect effects, with definition and detailed discussion deferred to the supplementary.

We construct both the Horvitz--Thompson estimator and Hajek-type estimator for all parameters. Details are in the supplementary.
      
\section{Further discussion}     
Our further efforts involve several aspects: (i) Develop asymptotic normality results for the point estimators and propose variance estimation strategy for inference. Our estimators fall broadly into weighting estimators with generalized propensity scores \citep{imbens2000role, aronow2017estimating}. We hope to further incorporate the between-cluster interference structure and provides a complete picture for statistical inference under such a setting. (ii) Propose a model-assisted approach for the estimation and inference. While we do not directly impose outcome models on our estimands, it is convenient to explore model-assisted approach to recover the point and variance estimators \citep{zhao2022reconciling}. (iii) Reanalyze the experiment by \cite{egger2022general} with the proposed interference framework. 

% 1. continuous exposure mapping
% 2. two bipartite networks
% 3. notation, subscript $ki$?


\bibliographystyle{apalike}
\bibliography{ref}



\newpage 
\appendix

\renewcommand{\thesection}{\Alph{section}}
\renewcommand{\thetheorem}{\Alph{section}.\arabic{theorem}}
\setcounter{theorem}{0}
\renewcommand{\thelemma}{\Alph{section}.\arabic{lemma}}
\setcounter{lemma}{0}
\renewcommand{\theproposition}{\Alph{section}.\arabic{proposition}}
\setcounter{proposition}{0}
\renewcommand{\thecorollary}{\Alph{section}.\arabic{corollary}}
\setcounter{corollary}{0}
\renewcommand{\thedefinition}{\Alph{section}.\arabic{definition}}
\setcounter{definition}{0}
\renewcommand{\thepage}{S\arabic{page}}
\setcounter{page}{1}
\renewcommand{\theequation}{\Alph{section}.\arabic{equation}}
\setcounter{equation}{0}
\renewcommand{\thetable}{S\arabic{table}}
\setcounter{table}{0}

\begin{center}
  \LARGE {\bf Supplementary Material}
\end{center}

\section{Between-cluster indirect effect}
For between-cluster indirect effects, similarly, we define the \textit{between-cluster controlled indirect effect} of $H_{i}$ on the outcome as
\begin{eqnarray*}
    \ie^{\textup{b}}(s,h,h^{\prime}) &=& n^{-1}\sumi Y_{i}(0,s,h) - n^{-1}\sumi Y_{i}(0,s,h^{\prime}),
\end{eqnarray*}
and the \textit{marginal between-cluster indirect effect} of $H_{i}$ on the outcome as
\begin{eqnarray*}
    \ie^{\textup{b}}(h,h^{\prime}) &=& n^{-1}\sumi E_{\bA_{(-i)}\mid H_{i} = h}\left\{Y_{i}(0,S_{i},h)\right\} - n^{-1}\sumi E_{\bA_{(-i)}\mid H_{i} = h^{\prime}}\left\{Y_{i}(0,S_{i},h^{\prime})\right\},
\end{eqnarray*}
where $E_{\bA_{(-i)}\mid H_{i} = h}(\cdot)$ denotes expectation conditional on $H_{i}=h$. The quantity $\ie^{\textup{b}}(s,h,h^{\prime})$ captures the effect of changing the level of between-cluster exposure from $h$ to $h^{\prime}$, with the own treatment fixed at zero and the within-cluster exposure $S_{i}$ fixed at level $s$. The quantity $\ie^{\textup{b}}(h,h^{\prime})$ is the corresponding marginal effect, depending on the distribution of treatments of other villages under the condition $H_{i}=h$ and $H_{i}=h^{\prime}$.

% \begin{remark}
%     When the treatment assignments $A_{i}$'s are independent, e.g., if we have Bernoulli randomization within each sublocation, then the conditional expectations $E_{\bA_{(-i)}\mid S_{i} = s}(\cdot)$ and $E_{\bA_{(-i)}\mid H_{i} = h}(\cdot)$ are equivalent to expectations over the marginal distribution of $H_{i}$ and $S_{i}$, respectively.
% \end{remark}

% \begin{remark}
%     All quantities are defined conditional on the realization of the first-stage randomization $\bZ$, which affects the distribution of $\bA$.
% \end{remark}

\section{Estimation}
We first construct an estimator for the village average potential outcome,
\begin{eqnarray*}
    \bar{Y}(a,s,h) &=& n^{-1}\sumi Y_{i}(a,s,h),
\end{eqnarray*}
for given $(a,s,h)$. Consider the Horvitz--Thompson estimator
\begin{eqnarray*}
    \hat{Y}^{\horthom}(a,s,h) &=& \frac{\sumi n^{-1}\indicator\{A_{i}=a, S_{i}=s, H_{i}=h\} Y_{i}}{\pr(A_{i}=a, S_{i}=s, H_{i}=h)},
\end{eqnarray*}
and the Hajek-type estimator
\begin{eqnarray*}
    \hat{Y}^{\haj}(a,s,h) &=& \frac{\sumi \indicator\{A_{i}=a, S_{i}=s, H_{i}=h\} Y_{i}}{\sumi \indicator\{A_{i}=a, S_{i}=s, H_{i}=h\}}.
\end{eqnarray*}

We then propose the following estimators for the direct and indirect effects defined in Section~\ref{sec::parameters}. For $*\in\{\horthom,\haj\}$, construct
\begin{eqnarray*}
    \hat{\de}^{*}(s,h) &=& \hat{Y}^{*}(1,s,h) - \hat{Y}^{*}(0,s,h), \\
    \hat{\ie}^{\textup{w},*}(s,s^{\prime},h) &=& \hat{Y}^{*}(0,s,h) - \hat{Y}^{*}(0,s^{\prime},h), \\
    \hat{\ie}^{\textup{b},*}(s,h,h^{\prime}) &=& \hat{Y}^{*}(0,s,h) - \hat{Y}^{*}(0,s,h^{\prime}).
\end{eqnarray*}

Under the randomization, the estimator $\hat{Y}^{\horthom}(a,s,h)$ is unbiased to $\bar{Y}(a,s,h)$. The Hajek estimator $\hat{Y}^{\haj}(a,s,h)$ is not unbiased in finite sample but is consistent to $\hat{Y}^{\haj}(a,s,h)$ and generally have more stable finite sample performance. Therefore, the corresponding Horvitz--Thompson estimators are unbiased, and the Hajek estimators are consistent.

We also construct estimators for the aggregated versions of the village average potential outcome, aggregated at various levels. Define
\begin{eqnarray*}
    \hat{Y}^{\horthom}(a) &=& \frac{\sumi\indicator\{A_{i}=a\} Y_{i}}{\sumi\pr(A_{i}=a)}, \quad \hat{Y}^{\haj}(a) \ =\ \frac{\sumi\indicator\{A_{i}=a\} Y_{i}}{\sumi\indicator\{A_{i}=a\}}, \\
    \hat{Y}^{\horthom}(a,s) &=& \frac{\sumi\indicator\{A_{i}=a, S_{i}=s\} Y_{i}}{\sumi\pr(A_{i}=a, S_{i}=s)}, \quad \hat{Y}^{\haj}(a,s) \ =\  \frac{\sumi\indicator\{A_{i}=a, S_{i}=s\} Y_{i}}{\sumi\indicator\{A_{i}=a, S_{i}=s\}}, \\
    \hat{Y}^{\horthom}(a,h) &=& \frac{\sumi\indicator\{A_{i}=a, H_{i}=h\} Y_{i}}{\sumi\pr(A_{i}=a, H_{i}=h)}, \quad \hat{Y}^{\haj}(a,h) \ =\  \frac{\sumi\indicator\{A_{i}=a, H_{i}=h\} Y_{i}}{\sumi\indicator\{A_{i}=a, H_{i}=h\}}.
\end{eqnarray*}
For $*\in\{\horthom,\haj\}$, estimators for the marginal quantities are therefore
\begin{eqnarray*}
    \hat{\de}^{*} &=& \hat{Y}^{*}(1) - \hat{Y}^{*}(0), \\
    \hat{\ie}^{\textup{w,*}}(s,s^{\prime}) &=& \hat{Y}^{*}(0,s) - \hat{Y}^{*}(0,s^{\prime}), \\
    \hat{\ie}^{\textup{b,*}}(h,h^{\prime}) &=& \hat{Y}^{*}(0,h) - \hat{Y}^{*}(0,h^{\prime}).
\end{eqnarray*}
Similarly, the corresponding Horvitz--Thompson estimators are unbiased, and the Hajek estimators are consistent.

      


\end{document}


% Here, $S_{i}$ indicates whether more than half of the villages are getting treated within the same sublocation as village $i$, while $H_{i}$ measures whether more half villages are treated among those connected to village $i$ but located in different sublocations. In other words, $S_{i}$ captures within-sublocation interference, since the potential outcome of village $i$ may depend not only on its own treatment, but also on the treatments of other villages in the same sublocation. On the other hand, $H_{i}$ captures between-sublocation interference, reflecting the possibility that nearby sublocations influence each other through spillover effects \citep{egger2022general}. 